% Part: applied-modal-logics
% Chapter: epistemic-logic
% Section: truth-at-w

\documentclass[../../../include/open-logic-section]{subfiles}

\begin{document}

\olfileid{aml}{el}{trw}

\olsection{په يوه نړۍ کښې رښتياوالی}

لکه د نورمال مودالي منطق په شان، هر معرفتي مدل ټاکي چې کوم !!{formula}s
ئې په کومو نړۍو کښې رښتيا ګڼل کېږي۔ د اړيکيزې معنٰی پوهنې د بنسټيز
مفهوم لپاره هماغه نښه کاروو: «مدل~$\mModel{M}$ په نړۍ~$w$ کښې
!!{formula}~$!A$ رښتيا کوي»۔ دا اړيکه په استقرايي ډول تعريفېږي، او د
ټولو غيرمودالي عاملونو لپاره ئې تعريف د نورمال مودالي حالت په شان عين دے۔

\begin{defn}\ollabel{defn:mmodels}
  په~$\mModel M$ کښې \emph{په~$w$ د !!a{formula}~$!A$ رښتياوالی}،
  چې په نښو کښې ئې $\mSat{M}{!A}[w]$ ليکو، په استقرايي ډول داسې تعريفېږي:
  \begin{enumerate}
  \tagitem{prvFalse}{\indcase{!A}{\lfalse}{هېڅکله $\mSat{M}{\lfalse}[w]$ نۀ وي}.}{}
  \tagitem{prvTrue}{\indcase{!A}{\ltrue}{تل $\mSat{M}{\ltrue}[w]$ وي}.}{}
  \item $\mSat{M}{p}[w]$ هله او يوازې هله چې $w \in V(p)$ وي۔
  \tagitem{prvNot}{\indcase{!A}{\lnot !B}{$\mSat{M}{\indfrm}[w]$ هله او يوازې هله چې
    $\mSat/{M}{!B}[w]$ وي}.}{}
  \tagitem{prvAnd}{\indcase{!A}{(!B \land !C)}{$\mSat{M}{\indfrm}[w]$ هله او يوازې هله چې
    $\mSat{M}{!B}[w]$ او $\mSat{M}{!C}[w]$ دواړه وي}.}{}
  \tagitem{prvOr}{\indcase{!A}{(!B \lor !C)}{$\mSat{M}{\indfrm}[w]$ هله او يوازې هله چې
    $\mSat{M}{!B}[w]$ يا $\mSat{M}{!C}[w]$ وي} (يا دواړه)۔}{}
  \tagitem{prvIf}{\indcase{!A}{(!B \lif !C)}{$\mSat{M}{\indfrm}[w]$ هله او يوازې هله چې
    $\mSat/{M}{!B}[w]$ يا $\mSat{M}{!C}[w]$ وي}.}{}
  \tagitem{prvIff}{\indcase{!A}{(!B \liff !C)}{$\mSat{M}{\indfrm}[w]$ هله او يوازې هله چې
    يا دواړه $\mSat{M}{!B}[w]$ او $\mSat{M}{!C}[w]$ وي، يا
    نۀ $\mSat{M}{!B}[w]$ وي او نۀ $\mSat{M}{!C}[w]$ وي}.}{}
  \item\ollabel{defn:sub:mmodels-box}
    \indcase{!A}{\Knows_a !B}{$\mSat{M}{\indfrm}[w]$ هله او يوازې هله چې
    د هر هغه $w' \in W$ لپاره $\mSat{M}{!B}[w']$ وي چې $R_a ww'$ وي}
  \end{enumerate}
\end{defn}

خو دلته د لاسرسي پر اړيکو د قيدونو په اړه فکر کول پکار دي۔ د
\olref{defn:sub:mmodels-box} د مادې له مخې، که هېڅ داسې~$w'$
نۀ وي چې $R_a ww'$ وي، نو !!a{formula}~$\Knows_a !B$ په~$w$
کښې رښتيا دے۔ دا د نورمال مودالي منطق هماغه ماده ده: که يوه نړۍ
جانشينې نړۍ ونۀ لري، نو ټول $\Box$-فارمولونه پکښې په تش ډول رښتيا
وي۔ که $\Knows$ د \emph{پوهې} د ښودلو عامل وګڼو، دا ډېره ناعادي
پايله ښکاري۔ ځکه موږ عموماً فکر کوو چې داسې \emph{هېڅ} حالت نشته
چې يو عامل په يوه وخت کښې هم $!A$ او هم $\lnot !A$ وپېژني۔

يوه لاره دا ده چې په معرفتي منطق کښې د لاسرسي اړيکه تل
\emph{انعکاسي} وټاکو۔ دا په نژدې معنا له دې نظر سره سمون لري چې
واقعي نړۍ د عامل له معلوماتو سره سازګاره ده۔ په حقيقت کښې معرفتي
منطقونه عموماً S5 کاروي، خو نور يې ښايي کمزوري نظامونه وکاروي،
چې دا د $\Knows_a$ د اړيکې له مطلوب تفسير سره تړاو لري۔

\begin{figure}
  \begin{center}
    \begin{tikzpicture}[modal]
      \node[world] (w1) [label={[align=right]right:\mTrue{p}\\ \mFalse{q}}]{$w_1$}; 
      \node[world] (w2) [label={[align=right]right:\mTrue{p}\\ \mTrue{q}},
        above right=of w1]{$w_2$}; 
      \node[world] (w3) [label={[align=right]right:\mFalse{p}\\ \mFalse{q}},
        below right=of w1] {$w_3$};
      \draw[<->] (w1) to node {$a$} (w2);
      \draw[->,loop left] (w1) to node {$a,b$} (w1);
      \draw[<->] (w1) to [swap] node {$b$} (w3);
      \draw[->,loop above] (w2) to node {$a,b$} (w2) ;
      \draw[->,loop below] (w3) to node {$a,b$} (w3) ;
    \end{tikzpicture}
  \end{center}
  \caption{يو ساده معرفتي مدل۔}
  \ollabel{fig:simple}
\end{figure}

\begin{prob}
  په \olref[aml][el][trw]{fig:simple} کښې وګورئ چې له لاندې کومې
  دعوې سمې دي:
  \begin{enumerate}
    \item $\mSat{M}{\lnot q}[w_1]$;
    \item $\mSat{M}{\Knows_a \lnot q}[w_1]$;
    \item $\mSat{M}{\Knows_b \lnot q}[w_1]$;
    \item $\mSat{M}{\Knows_b q \lor \Knows_b \lnot q}[w_2]$;
    \item $\mSat{M}{\Knows_a( \Knows_b q \lor \Knows_b \lnot q)}[w_2]$;
    \item $\mSat{M}{\EKnows_{\{a,b\}} \lnot q}[w_3]$;
    \end{enumerate}
\end{prob}

اوس چې په يوه نړۍ کښې مو د رښتياوالي بنسټيز تعريف ورکړ، د نورمال
مودالي منطق نور معنايي مفاهيم، لکه مودالي اعتبار او لازم راتلل، هم
دغه نوي تفسير ته د مودالي عاملونو په پلي کولو سره راځي۔

اوس د مشترکې پوهې د عامل~$\CKnows_G$ لپاره هم د صدق شرطونه ټاکلے
شو۔ د \olref[sfr][rel][ops]{sec} له مخې ياد کړئ چې د اړيکې~$R$
\emph{انتقالي تړون}~$R^+$ داسې تعريفېږي:
\begin{align*}
  R^+ &= \bigcup_{ n \in \mathbb{N}} R^n, 
\intertext{چې}
  R^0 & = R \text{ او}\\ 
  R^{n+1} & = \Setabs{\tuple{x, z}}{\exists y (R^n xy \land Ryz)}.
\end{align*}
بيا، که $G$ د عاملانو يوه ډله وي، نو $R_G = ( \bigcup_{b
\in G} R_b )^+$ د ټولو عاملانو د لاسرسي اړيکو د اتحاد انتقالي
تړون تعريفوو۔

\begin{defn}
که $G' \subseteq G$ وي، نو $\mSat{M}{\CKnows_{G'} !A}[ w]$
هله او يوازې هله وي چې د هر هغه $w'$ لپاره چې $R_{G'} w w'$
وي، $\mSat{M}{!A}[w']$ وي۔
\end{defn}

\end{document}
